\documentclass[11pt]{article}
\title{Temporal Coupling in Xenopus Tail Regeneration: An Identifiability Limit and a Testable Timing-Gate Hypothesis}
\date{}
\author{SIGMA AI-assisted research synthesis}
\usepackage[T1]{fontenc}
\usepackage[utf8]{inputenc}
\usepackage[margin=1in]{geometry}
\usepackage{amsmath,amssymb,array,booktabs,longtable,pdflscape}
\usepackage{microtype}
\usepackage[hyphens]{url}
\usepackage{titlesec}
\emergencystretch=2em
\raggedbottom
\linespread{1.07}
\usepackage{xcolor}
\definecolor{blue}{HTML}{24485E}
\titleformat{\section}{\Large\bfseries\sffamily\color{blue}}{\thesection}{0.8em}{}
\titleformat{\subsection}{\large\bfseries\sffamily}{\thesubsection}{0.7em}{}
\setlength{\parskip}{.6em}
\setlength{\parindent}{0pt}
\begin{document}
\maketitle
\begin{center}\small\textit{AI-generated research synthesis \textperiodcentered{} SIGMA}\end{center}

\begin{abstract}
Why does tail-regeneration competence decline temporarily during Xenopus laevis development, and what measurement could distinguish failed injury signaling from failed cellular responsiveness? This paper develops a bounded analytical answer rather than reporting a discovered mechanism. If early injury measurements and later regenerative-state measurements come from unlinked animals, their stage-specific distributions cannot identify individual temporal coupling, even its sign. A constructive counterexample establishes this limitation; whether it applies to the nominated public datasets remains unresolved because their sampling structure has not been verified.

A prospective hypothesis instead distinguishes insufficient early epidermal signaling from delayed epidermal receptivity through equal-exposure pulses delivered at different wound ages. The restricted models predict opposite transition-drive contrasts, but molecular readouts must separate wound age from response latency. No biological effect estimate, predictive improvement, novelty determination, or independent validation is available. The contribution is an explicit boundary on retrospective inference and a falsifiable experimental question, not an explanation already established for developmental loss and recovery of regeneration.
\end{abstract}
\section*{About this paper}
This paper was generated by SIGMA, an AI-assisted research system, from the project's recorded investigations, findings, and supporting evidence.
The analysis and conclusions reflect an automated synthesis of that research. Although SIGMA incorporates evidence tracking, critical review, and consistency checks, the resulting paper has not necessarily undergone independent verification or peer review.
Readers should evaluate its arguments, supporting evidence, and stated limitations before relying on its conclusions.
Research style: Distinguish observed biology from mathematical identifiability, mechanistic hypotheses and proposed interventions. Identify organism and developmental stage, biological unit, whether specimens are paired, sample and assay provenance, and cross-species limits. Compare mechanisms through predictions that would genuinely discriminate them; state matched-time controls, alternative explanations and falsifiers without inventing results.
\section*{Evidence coverage limitation}
The evidence index omitted 3390 material saved research field fragment(s) across 10 source(s). These fragments were not reviewed for this paper. Interpretations that could depend on omitted material remain provisional; the complete frozen source snapshot and omission ledger are retained with the generation record.
\textit{Claims discussed: }C11
\section{Introduction and biological context}
The motivating question concerns developmental changes in tail-regeneration competence in Xenopus laevis, including variation between animals. The proposed target is not another association with regeneration-initiating cells, extracellular-matrix remodeling, or familiar injury pathways, but a stage-dependent relationship between an early basal-epidermal signal and a later regenerative transition. A comparison of stage classes alone cannot establish that relationship within animals or explain subsequent recovery of competence.

The supplied biological background treats basal-epidermal origins of regeneration-initiating cells and regeneration-associated matrix-remodeling defects as prior knowledge. Those exclusions define the proposed contribution, but the cited articles and subsequent literature have not been inspected sufficiently to establish originality. Evidence concerning Xenopus tropicalis is not transferred to Xenopus laevis as support for its developmental trajectory or a nutritional mechanism.

The immediate answerable question is whether the proposed temporal relationship is identifiable from the available sampling design and, if not, what intervention could distinguish two restricted mechanisms. The paper therefore separates a conditional mathematical result from a prospective biological hypothesis. It does not claim that either constitutes a completed transcriptomic reanalysis or a newly validated regeneration mechanism.

\textit{Claims discussed: }C11, C12
\section{Methods and evidence}
The analysis synthesized the supplied source-linked excerpts and distinguished reported background, explicit assumptions, conditional deductions, conjectured mechanisms, and outstanding evidence requirements. Dataset descriptions were treated as provisional accession-level pointers rather than inspected sample manifests. No expression matrices, specimen-level outcomes, or original sample records were analyzed.

The analytical component examined an animal-level coupling estimand, constructed observationally equivalent models with different coupling coefficients, and considered the rank deficiency created by assigning a stage-average early response to every later specimen. These arguments address what information a sampling design can identify; they do not estimate a biological association or establish that the nominated accessions have that design.

The mechanistic component compared two restricted signal-receptivity models under matched integrated epidermal exposure. Their conditional predictions were evaluated alongside the saved timing objection, alternative downstream blocks, assay-feasibility requirements, and proposed validation criteria. The resulting experimental outline is prospective: no pulse intervention, reporter validation, predictive benchmark, or independent biological replication was performed.

\textit{Claims discussed: }C10, C11
\section{Biological question and observational scope}
The biological target is a developmental change in the relationship between early injury signaling and subsequent entry into a regenerative epidermal state in Xenopus laevis tadpoles, rather than a molecular signature that merely separates developmental-stage labels. The nominated comparison is between stages $40$ and $46$, described in the supplied study framing as regeneration-competent and refractory conditions, respectively. Because competence varies among animals, the proposed relationship asks whether an animal\textquoteright{}s early wound response predicts its later basal-epidermal regeneration-initiating-cell and extracellular-matrix response, and whether that relationship differs between stages.

This is a candidate dependency to investigate, not an observed biological result.

Regeneration-initiating cells (RICs) and extracellular-matrix (ECM) remodeling delimit the question rather than constitute its proposed novelty. The supplied biological account attributes transient RIC formation from basal epidermal cells and ECM-remodeling defects associated with regeneration failure to the 2024 Genome Biology study identified by DOI 10.1186/s13059-024-03396-3. It likewise treats reactive oxygen species and inflammatory and developmental signaling as prior knowledge, so their detection alone would not establish the proposed dependency.

These are supplied prior-knowledge exclusions: the underlying publications and subsequent literature have not been comprehensively examined, and no priority claim is made for the candidate relationship.

The observational starting point is GEO GSE245320 and its nominated component series: bulk RNA sequencing in GSE245317 and GSE245318, single-cell RNA sequencing in GSE245312, and spatial profiling in GSE245313. The supplied description assigns the bulk comparison measurements at $0$, $0.5$, and $6$ hours post-amputation and at $3$ days. That schedule identifies a potentially useful separation between early injury responses and later regenerative-state measurements, but it is not a verified sample-by-condition design. These accessions are described as resources from a shared study and therefore cannot be counted as independent biological replications merely because they have separate accession numbers.

The assays have different potential roles in addressing the question. Bulk measurements could support stage-dependent population expression contrasts if replication and design structure permit them, whereas the predominantly regenerating single-cell material could help define candidate epidermal states and response modules. Such single-cell material is not a matched survey of both regeneration conditions and cannot substitute for missing refractory observations or establish early-to-late identity within an animal. The spatial accession is also a nominated resource, not evidence that the required stages, wound ages, tissue boundaries, or linked specimens are available.

Actual sampling provenance determines the biological interpretation of any comparison. Original sample tables and expression files were not available in the supplied evidence, leaving animal versus pool identity, pool size, early-to-late pairing, replicate counts, clutch and batch structure, tissue boundaries, phenotype labels, and assay availability unresolved. The relevant sampling unit must therefore be established before analysis: technical libraries or cells from one specimen are not additional biological replicates, while an independently collected pool may support pool-level inference without representing an individual animal.

Species and developmental stage at both amputation and collection also require explicit resolution; nutritional and stage-dependent findings from X. tropicalis cannot be imported as equivalent X. laevis conditions.

This scope makes the distinction between a population trajectory and an animal-level dependency central to the argument. A stage-associated change in early or late expression would address the nominated developmental contrast, but the proposed temporal relationship additionally requires evidence about how those measurements are linked across biological units. The available material therefore sharpens the question from whether competent and refractory samples differ to whether their sampling structure can support a test of early-signal-to-regenerative-state coupling.

\section{Why unlinked snapshots cannot identify animal-level coupling}
Unlinked early and late measurements cannot, by themselves, establish whether developmental stage changes the relationship between an individual animal\textquoteright{}s injury response and its subsequent regenerative state. For the Xenopus laevis tadpole question, let $S$ denote developmental stage, $E$ an animal\textquoteright{}s early wound-response measurement, and $L$ its later basal-epidermal regeneration-initiating-cell and extracellular-matrix transition measurement. A candidate measure of temporal coupling is the within-stage linear-projection slope $\beta_s=\operatorname{Cov}(E,L\mid S=s)/\operatorname{Var}(E\mid S=s)$, provided the relevant moments exist and $\operatorname{Var}(E\mid S=s)>0$.

This estimand concerns variation among animals within a stage, rather than differences between stage averages; even when identifiable, it is an association rather than a causal effect.

The precise premise of the limitation is that the available observations identify only the stage-conditioned marginal distributions $P(E\mid S=s)$ and $P(L\mid S=s)$, without animal linkage or justified restrictions that identify their joint distribution $P(E,L\mid S=s)$. Independent destructive collections at early and late times can satisfy this premise: matching developmental stage, batch, or clutch does not make different animals repeated measurements of one individual. Whether this premise describes GSE245320 and its nominated component series remains unresolved, because the supplied evidence does not establish their actual animal identities, pooling, or early-to-late linkage.

The argument therefore establishes a conditional mathematical boundary, not a verified description of those datasets.

A two-point construction demonstrates why even the sign of coupling is not generally identifiable. Within stage $s$, let $Z$ take values $-1$ and $+1$ with equal probability, and set $E=\mu_s+Z$. In a same-sign model, set $L=\nu_s+Z$; in an opposite-sign model, set $L=\nu_s-Z$. Both models give exactly the same early marginal distribution, with values $\mu_s-1$ and $\mu_s+1$, and the same late marginal distribution, with values $\nu_s-1$ and $\nu_s+1$, each occurring equally often. These are constructed mathematical values, not measured or simulated tadpole expression data.

The difference lies entirely in the unobserved correspondence between early and late values. In both models, $\operatorname{Var}(E\mid S=s)=1$, but the same-sign model has $\operatorname{Cov}(E,L\mid S=s)=+1$ and $\beta_s=+1$, whereas the opposite-sign model has covariance and slope $-1$. Consequently, the distributions of independently sampled, unlinked early and late observations are identical under models with opposing animal-level relationships. Perfect knowledge of both marginals would not resolve that ambiguity.

The construction also prevents identification of the direction of a developmental-stage interaction. Assigning the same-sign model to stage $40$ and the opposite-sign model to stage $46$ gives $\beta_{46}-\beta_{40}=-2$; reversing the assignments gives $\beta_{46}-\beta_{40}=+2$. Every stage-specific early and late marginal remains unchanged between these assignments, including its mean and variance. Thus stage-specific differences in average expression can coexist with mutually incompatible within-animal coupling contrasts: observing the former does not determine the latter.

Replacing missing individual early measurements with stage-level early means creates a related algebraic failure. For this calculation, recode the two stages as $S\in\{0,1\}$ and assign each late sample the early mean of its stage, $\widetilde E=e_0+(e_1-e_0)S$. Because $S\widetilde E=e_1S$, the columns for an intercept, stage, assigned early response, and their interaction all lie in the span of the intercept and stage; the design matrix has rank at most $2$. The proposed early-response interaction is therefore not separately identifiable, regardless of how many late biological replicates receive those same assigned means.

Additional unlinked animals can improve precision for marginal distributions, and additional unlinked time points can characterize population trajectories, but neither supplies the missing individual correspondence without further identifying information. Correlations among stage or time-group means concern an ecological relationship, not $\beta_s$. Likewise, linked pool measurements could support a pool-level analysis without automatically identifying animal-level coupling, and additional technical libraries from one specimen do not create independent biological replicates.

This boundary does not rule out descriptive differential expression, appropriately identified stage-by-time contrasts, or prediction from measurements genuinely available on an individual before its outcome. Actual longitudinal linkage, informative interventions, or justified structural restrictions could change what is identifiable. The consequence for the central argument is therefore specific: a stage-dependent population trajectory cannot establish that strong early responders subsequently enter a regenerative state more or less effectively. That claim requires information about the early-to-late relationship itself, not merely more precise snapshots of its separate components.

\section{A restricted timing-gate hypothesis}
A restricted timing-gate model makes the distinction between insufficient injury signaling and delayed cellular responsiveness experimentally explicit. For an individual Xenopus laevis tadpole at developmental stage $s$, let $E_s(t)$ denote basal-epidermal ERK activity at time $t$ after amputation, and let $C_s(t)$ denote receptivity to enter a transient regenerative epidermal state. Define cumulative transition drive as $J_s=k\int E_s(t)C_s(t)\,dt$, with $k>0$, over an interval containing the candidate receptive windows. This is a prospective model, not a fit to observed animal trajectories: receptivity is a latent functional property, and transition drive is neither a measured state score nor an established measure of tail regrowth.

The product in this model expresses a specific conditional dependency: increasing epidermal activity contributes to transition drive only while the epidermis is receptive. The insufficient-early-signal model, $A_{\mathrm{amp}}$, attributes failure at stage $46$ to inadequate early ERK activity while retaining early receptivity at both stages $40$ and $46$. The delayed-receptivity model, $B_{\mathrm{delay}}$, instead retains early receptivity at stage $40$ but places the receptive window at stage $46$ later after injury. In the latter hypothesis, basal epidermal cells remain inducible, but their receptive period may occur after the endogenous signal has declined.

These alternatives are deliberately narrower than a general failure to induce regeneration-initiating cells or extracellular-matrix remodeling. Either inadequate activity during an intact receptive window or displacement of that window could reduce entry into the regenerative epidermal state. Consequently, observing reduced state induction would not distinguish the two models. Their distinguishing prediction concerns the timing of an added signal, rather than the presence or absence of the downstream state alone.

Let $p_{\mathrm{early}}(t)$ and $p_{\mathrm{late}}(t)$ be nonnegative additions to basal-epidermal ERK activity, supported in the respective early and late windows, with equal positive integrated activity: $\int p_{\mathrm{early}}(t)\,dt=\int p_{\mathrm{late}}(t)\,dt=Q>0$. Within stage $46$, hold the untreated activity trajectory, receptivity trajectory, and other background conditions fixed between the two counterfactual treatments. Substitution into the model gives $J_{46,\mathrm{late}}-J_{46,\mathrm{early}}=k\int C_{46}(t)[p_{\mathrm{late}}(t)-p_{\mathrm{early}}(t)]\,dt$, because the untreated contribution cancels.

Thus, matched integrated activity isolates a pulse-weighted difference in receptivity under these assumptions, rather than a difference in total added signal.

The opposite signs follow exactly for idealized, nonoverlapping receptive windows. Under $A_{\mathrm{amp}}$, suppose $C_{46}(t)=c>0$ throughout the early pulse window and $C_{46}(t)=0$ throughout the late window; then $J_{46,\mathrm{late}}-J_{46,\mathrm{early}}=-kcQ<0$. Under $B_{\mathrm{delay}}$, reverse those supports while retaining the same positive receptive level, giving $J_{46,\mathrm{late}}-J_{46,\mathrm{early}}=+kcQ>0$. Early supplementation therefore produces greater transition drive in the restricted amplitude model, whereas late supplementation produces greater drive in the restricted delay model, without requiring unequal total added ERK activity.

For gradual or overlapping receptive windows, the integral contrast remains applicable, but neither mechanism's name guarantees its sign. Late supplementation is favored only when its activity overlaps greater receptivity than does the early supplementation; equal pulse-weighted receptivity gives no contrast. Both restricted models retain early receptivity at stage $40$, so their disagreement is specifically about the stage-$46$ response to pulse timing. The discriminating prediction is therefore not simply that ERK supplementation can help, but that changing its delivery time reverses the preferred intervention between the stipulated stage-$46$ mechanisms.

This interpretation requires the pulses to represent comparable achieved ERK perturbations in basal epidermis, rather than merely equal treatment durations or nominal doses. It also assumes that the intervention does not independently shift receptivity or introduce timing-dependent injury or growth effects. Cell-specific delivery, exposure matching, and a sufficiently sensitive state-transition assay remain unvalidated in the supplied evidence. Moreover, the contrast concerns cumulative transition drive: a single molecular endpoint can still reflect unequal response latency or decay rather than a wound-age-dependent receptive window.

Transferring the predicted sign to regrowth requires an additional monotone, nonsaturated downstream relationship without a separate dominant block. Saturation, absent receptivity in both tested windows, or a downstream obstruction could erase a regrowth contrast even when the proposed signal manipulation is achieved. The analytical result therefore sharpens the biological question from whether regenerative-state induction fails to whether equal added epidermal activity is more effective early or late\textemdash{}a distinction between two specified mechanisms, not between every possible signaling defect and every possible state-induction defect.

\section{Discriminating timing from kinetics and alternative gates}
A pulse-timing experiment can distinguish the restricted signaling models only if it separates wound-age-dependent responsiveness from the kinetics of the measured response. In proposed experiments on Xenopus laevis tadpoles at stages $40$ and $46$, the relevant comparison is therefore not simply whether late treatment produces more regenerative-state signal at one endpoint. It is whether matched basal-epidermal signaling exposures produce different vehicle-adjusted induction trajectories when delivered at different wound ages, and whether those trajectories precede remodeling and regrowth. Temporal association or rescue alone would leave downstream competence and alternative gates unresolved [C10].

The intervention should be calibrated against achieved epidermal ERK activity rather than nominal treatment duration. A basal-epidermal light-gated RAF/ERK actuator is one proposed implementation, but its availability, specificity, reversibility, and noninjurious operation in these cells remain unestablished; temporally resolved state measurements likewise require validation. Separate calibration specimens should establish comparable pulse amplitude and integrated added activity within a prespecified tolerance. If early endogenous activity is already near saturation, increasing illumination does not necessarily create an exposure comparable to a late pulse, and an unmatched comparison cannot be interpreted as a timing effect.

A fixed molecular endpoint confounds receptivity with elapsed time since stimulation. At $6$ hours post-amputation, for example, an early pulse has had longer than a late pulse to induce a transient state and allow that state to decay. A larger late-pulse response could consequently reflect more recent induction rather than greater cellular responsiveness at the later wound age. Stage-dependent induction or decay kinetics could also generate a stage-by-pulse-timing interaction without the proposed shift in receptivity.

The molecular design should therefore combine measurements at common wound ages with measurements at matched post-pulse latencies. Let $t_p$ denote the end of pulse $p$, and let $\tau=t-t_p$ denote elapsed time after that pulse. For a prespecified basal-epidermal state readout $R$, define the proposed vehicle-adjusted trajectory as $\Delta_{s,p}(t)=\mathbb{E}[R(t)\mid s,p]-\mathbb{E}[R(t)\mid s,\mathrm{vehicle}_p]$, where vehicle specimens receive handling matched to the delivery schedule. Comparing these trajectories at a common $t$ addresses state occupancy at the same wound age; comparing them at common values of $\tau$ addresses response development after equivalent exposure. Neither comparison substitutes for the other.

Serial measurements must cover induction and subsequent persistence rather than merely add a second isolated endpoint. If early- and late-pulse trajectories become comparable after alignment by post-pulse latency, a difference at the common endpoint is consistent with response kinetics and does not identify a wound-age-dependent gate. A late-pulse advantage that persists across matched latencies would be more consistent with greater late responsiveness, provided exposure, cell identity, and background conditions are comparable. Even then, developmental progression and metabolic state remain potential determinants of that responsiveness rather than excluded explanations.

The response assay must distinguish entry into the proposed regenerative epidermal state from a change in the amount of epidermis sampled. Prespecified per-basal-cell state measurements should therefore be accompanied by absolute epidermal abundance and subsequent ECM-remodeling measurements, with stable normalization across stages and treatment times. Animals should be randomized within developmental stage and clutch, with cells and image fields treated as subsamples rather than independent biological replicates.

A validated nonlethal reporter could link signaling, state induction, and later regrowth within an individual; destructive molecular companion cohorts would instead provide unpaired group trajectories, not individual mediation measurements.

Reciprocal inhibition and downstream bypass address different parts of the causal interpretation. Under both restricted models, early inhibition at stage $40$ is predicted to impair the transition more than late inhibition, so that comparison is a necessity control rather than the decisive test between the models. A reversible MEK inhibitor is a candidate implementation, but systemic inhibition would establish neither epidermal specificity nor a basal-cell site of necessity without compartment-resolved target-engagement evidence.

A separately validated downstream bypass should accompany upstream stimulation to ask whether the regenerative state remains accessible when early signaling is circumvented; no suitable bypass is established by the supplied evidence.

The alternative explanations in the table require controls beyond the pulse comparison itself. Wounded vehicle controls, sham or unamputated controls, actuator-negative illumination controls, and actuator-positive no-light controls should distinguish signaling effects from delivery, injury, and developmental effects. Amputation plane, feeding, density, temperature, observation time, and processing batches should be standardized or balanced, while closure, viability, and tissue composition are measured directly. A favored pulse must improve over vehicle: a late-versus-early advantage caused only by early-treatment damage is not rescue.

Rescue without a preceding detected state transition would challenge the specified mediation model only if the assay adequately covers the relevant cells, times, and detection range. Conversely, failure of both pulses despite target engagement would remain compatible with persistent downstream incompetence, another limiting gate, or incorrectly chosen windows; an imprecise null would not discriminate among them. These requirements change the proposed timing contrast from a putative mechanism identifier into a conditional test of responsiveness: evidence for a timing gate requires resolved molecular trajectories and causal controls, not merely a pulse-associated increase in regrowth [C10].

\begin{longtable}{@{}>{\raggedright\arraybackslash}p{\dimexpr(\linewidth-2\tabcolsep*2)/126*42\relax}>{\raggedright\arraybackslash}p{\dimexpr(\linewidth-2\tabcolsep*2)/126*42\relax}>{\raggedright\arraybackslash}p{\dimexpr(\linewidth-2\tabcolsep*2)/126*42\relax}@{}}
\caption{Alternative explanations for an apparent pulse-timing benefit and the proposed comparisons needed to distinguish them. These are conditional predictions and design requirements, not observed outcomes.}\\
\toprule
\textbf{Alternative explanation} & \textbf{Discriminating comparison or measurement} & \textbf{Interpretive boundary}\\
\midrule
\endfirsthead
\toprule
\textbf{Alternative explanation} & \textbf{Discriminating comparison or measurement} & \textbf{Interpretive boundary}\\
\midrule
\endhead
\midrule
\multicolumn{3}{r}{\footnotesize Continued on next page}\\
\endfoot
\bottomrule
\endlastfoot
Tissue-composition change & Measure the regenerative-state response within basal epidermal cells alongside absolute epidermal abundance. & A larger aggregate signal without a corresponding within-cell response is compatible with altered abundance rather than state induction.\\
Generic growth or wound effects & Compare each pulse with schedule-matched wounded vehicle controls; measure closure, viability, state induction, remodeling, and regrowth separately. & Improved regrowth alone does not establish mediation through the proposed epidermal transition.\\
Developmental progression or metabolic gating & Record developmental progression and metabolic state under standardized feeding and environmental conditions. & A timing benefit does not exclude these gates; standardization alone does not identify the molecular basis of responsiveness.\\
Persistent downstream incompetence & Combine verified upstream target engagement with a separately validated downstream-bypass test. & Failure of both upstream pulses is not specific to absent receptivity and can also reflect another downstream block or unsuitable windows.\\
Systemic rather than epidermal necessity & Measure epidermal and non-epidermal target engagement during reversible inhibition and washout. & A systemic blocking effect cannot by itself establish basal-epidermal necessity.\\
\end{longtable}
\par\smallskip{\footnotesize Evidence: [Research session 21, artifact 1], [Research session 21, artifact 4], [Research session 21, artifact 5]; claims: C10}
\textit{Claims discussed: }C10
\section{Validation thresholds and developmental recovery}
The timing-gate hypothesis requires distinct standards for mathematical discrimination, predictive association and biological mechanism; satisfying one does not establish the others. Numerical contrasts in an illustrative construction establish what follows from stipulated assumptions, not the size of a regenerative response in Xenopus laevis tadpoles. Likewise, a precise conditional model and a proposed minimum effect do not supply an observed effect estimate.

The supplied evidence contains no measured biological effect with an uncertainty interval, demonstrated predictive improvement or independent biological confirmation, so the proposed interaction remains a question for testing rather than a discovered explanation of the developmental switch [C11].

A reproducible computational association would first require an estimand supported by the actual biological sampling units. Animal-level prediction requires linked measurements and outcomes from individual animals; independently collected pools can support pool-level inference but cannot silently substitute for individuals. Evaluation splits and uncertainty estimates must respect those units, rather than count cells or technical libraries as independent specimens. Moreover, reuse of component datasets from the same study cannot provide the independent biological validation needed to extend an association beyond its original sampling context [C11].

Incremental predictive value should be assessed against preregistered stage-only, regeneration-initiating-cell/extracellular-matrix-only and known-signaling-module baselines, using a fixed primary metric and uncertainty procedure. Molecular definitions, outcome scoring and multiplicity handling should be fixed before confirmatory evaluation. The prediction target and measurement times also matter: a later molecular endpoint cannot serve as its own predictor, and a destructive state assay cannot furnish an individual baseline for an animal subsequently followed to regrowth.

Such a comparison would require a validated live surrogate or an explicitly narrower set of feasible baselines; separate cohorts would not restore the missing individual linkage [C11].

Quantitative decision thresholds must be justified independently of the confirmatory outcomes. A separate pilot should establish assay performance, estimate between-animal and between-clutch variability, and inform both a biologically meaningful minimum effect and a finite confirmatory sample size. Proposed regrowth contrasts, exposure tolerances and enrollment numbers are design choices, not empirical findings or evidence of adequate precision. Confirmatory reporting should therefore pair the estimated interaction or rescue effect with uncertainty and its prespecified interpretation, rather than treat passage of an arbitrary numerical threshold as validation [C11].

Negative results require the same separation between biological magnitude and statistical precision. A nonsignificant rescue comparison would not exclude a useful effect if its uncertainty remained broad; exclusion would require valid exposure and assay checks together with sufficient precision to rule out a prespecified meaningful benefit. Conversely, rescue without a detected regenerative-state transition would challenge the proposed mediation only if the state assay adequately covered the relevant basal epidermal cells and response times with established sensitivity. These criteria distinguish evidence against a restricted mechanism from an inconclusive experiment [C11].

Mechanistic support would additionally require independent specimens and preferably orthogonal confirmation of the proposed links from epidermal signal timing to state induction, remodeling and regrowth. An observational interaction alone cannot distinguish impaired signal transmission from intrinsic inability to enter the regenerative state, while pulse rescue alone does not identify the cellular decoder. Rescue, blockade and downstream-bypass evidence must therefore be interpreted together, with validated molecular readouts and the timing controls needed to distinguish receptivity from response kinetics. These are requirements for a causal interpretation, not results already supplied by the conditional model [C11].

Developmental recovery imposes a separate boundary on that interpretation. A comparison restricted to stages $40$ and $46$ cannot establish the later return of competence, and no later-competent X. laevis condition has been verified in the supplied evidence. Extending the hypothesis to recovery would require independently establishing competence in an appropriate later condition and testing the proposed relationship there, rather than importing a stage pattern from X. tropicalis [C12].

The argument is therefore narrowed in a consequential way: even successful discrimination of the restricted models at the two nominated stages would address developmental loss, whereas an explanation of temporary loss and subsequent recovery requires evidence that the same mechanism also accounts for the verified return of regenerative ability.

\textit{Claims discussed: }C11, C12
\section{Findings and hypotheses}
\subsection{C10 -- Unresolved}
A temporal association or pulse rescue would not by itself distinguish signal transmission from absent downstream competence, tissue-composition changes, metabolic gating, developmental progression, or generic wound effects.
The saved argument requires downstream-bypass and reciprocal-blocking tests alongside sham, vehicle, wound-size, closure, viability, and composition controls. Rescue without a detected regenerative-state change challenges mediation only if the assay adequately samples the relevant cells and times. These alternatives remain live explanations because no discriminating biological observations are supplied.
\textit{Evidence: }[Research session 21, artifact 1], [Research session 21, artifact 5], [Research session 21, artifact 4]
\subsection{C11 -- Limitation}
The supplied evidence does not establish a numerical biological effect, uncertainty interval, improved predictive performance, independent replication, or a new mechanism for the developmental switch.
The saved outcome explicitly reports missing empirical effects and independent validation, while accession applicability and novelty remain unassessed. Mathematical values in illustrative constructions are not gene-expression measurements or treatment effects. The supported outcome is a conditional inference limit and an unconfirmed experimental hypothesis, not a replicated computational discovery.
\textit{Evidence: }[Research session 21], [Research session 21, artifact 5]
\subsection{C12 -- Unresolved}
Explaining later recovery of regeneration competence requires a separately verified later-competent Xenopus laevis condition; a comparison restricted to stages $40$ and $46$ cannot establish the full developmental trajectory.
The saved obligations explicitly distinguish loss of competence from its later recovery. No later-competent condition has been verified or tested in the supplied evidence. Cross-species stage patterns cannot fill that gap.
\textit{Evidence: }[Research session 21, artifact 1], [Research session 21, artifact 4], [Research session 21, artifact 5]
\subsection{C1 -- Limitation}
The supplied evidence does not establish a numerical biological effect, uncertainty interval, improved predictive performance, independent replication, or a new mechanism for the developmental switch.
The saved outcome explicitly reports missing empirical effects and independent validation, while accession applicability and novelty remain unassessed. Mathematical values in illustrative constructions are not gene-expression measurements or treatment effects. The supported outcome is a conditional inference limit and an unconfirmed experimental hypothesis, not a replicated computational discovery.
\textit{Evidence: }[Research session 21], [Research session 21, artifact 5]
\subsection{C2 -- Unresolved}
Explaining later recovery of regeneration competence requires a separately verified later-competent Xenopus laevis condition; a comparison restricted to stages $40$ and $46$ cannot establish the full developmental trajectory.
The saved obligations explicitly distinguish loss of competence from its later recovery. No later-competent condition has been verified or tested in the supplied evidence. Cross-species stage patterns cannot fill that gap.
\textit{Evidence: }[Research session 21, artifact 1], [Research session 21, artifact 4], [Research session 21, artifact 5]
\section{Limitations}
\begin{itemize}
\item Original sample manifests, matrices, specimen outcomes, pooling, pairing, replicate counts, batches, expression units, and assay availability have not been verified; the conditional sampling argument cannot yet be attributed to the actual GEO series.
\item Prior-art exclusions come from the supplied project description. Original articles and later publications have not been inspected sufficiently to support a claim of novelty or an exhaustive literature account.
\item The pulse comparison distinguishes only two deliberately restricted models. Broad regenerative-state failure can include both, and translating transition drive into regeneration requires additional downstream assumptions.
\item A fixed post-amputation molecular endpoint confounds pulse timing with elapsed response time. Exposure matching, selective perturbation, reporter specificity, and an adequate sampling schedule remain feasibility requirements.
\item There are no biological effect estimates, uncertainty intervals, completed predictive benchmarks, independent biological validations, or verified observations of later regained competence in the supplied evidence.
\end{itemize}
\section{Open obligations}
\begin{itemize}
\item Acquire original GEO series and sample records for GSE245320, GSE245317, GSE245318, GSE245312, and GSE245313; establish a row-per-library manifest covering organism, exact stage and time, animal versus pool, biological replication, batch, units, accessibility, longitudinal linkage, and phenotype labels.
\item Determine which population contrasts are identifiable after accounting for stage and batch. Stop animal-level temporal-coupling analysis if only unlinked snapshots exist, and do not replace individual measurements with stage-average predictors.
\item Inspect the cited original publications and subsequent work through the investigation date; document the exact interaction or timing dependency tested before assigning originality.
\item Validate a selective, noninjurious basal-epidermal intervention and temporally resolved regenerative-state measurements; demonstrate that early and late exposures can be matched and that wound age can be distinguished from response latency.
\item Freeze molecular modules, homeolog handling, prediction targets, biological-unit splits, stage-only, RIC/ECM-only and known-signaling baselines, primary metrics, multiplicity handling, and clutch-aware uncertainty procedures before confirmatory analysis.
\item Use a separate pilot to justify a minimum-interest or equivalence margin and a finite confirmatory sample size. A nonsignificant or imprecise result must remain inconclusive rather than count as falsification.
\item Obtain independent randomized biological evidence with reciprocal blocking, downstream-bypass, wound-healing, viability, composition, developmental, and metabolic controls before claiming a mechanism.
\item Verify and test a later-competent Xenopus laevis condition before extending an early-versus-refractory comparison to an explanation of developmental recovery.
\end{itemize}
\section{Frozen obligation reconciliation}
Unlisted obligations are not presumed resolved. The archived research snapshot retains complete source references and the full reconciliation.
\begin{itemize}
\item ChatGPT plan or app usage limit reached. Manage usage at \url{https://chatgpt.com/settings/usage}; no API-key fallback occurred. (HTTP 400) No strategy was saved. \textit{ (Not listed by author; not presumed resolved; 1 saved source reference(s))}
\item No research artifact was saved for this strategy. \textit{ (Not listed by author; not presumed resolved; 3 saved source reference(s))}
\item Acquire original GEO sample metadata and files; verify species, exact stages/times, independent animals versus pools, pairing, batches, replicate counts, outcome labels and assay availability before analysis. If only unlinked snapshots exist, stop the animal-level temporal-coupling analysis. \textit{ (Not listed by author; not presumed resolved; 3 saved source reference(s))}
\item Inspect the cited original papers and later literature through the investigation date; record exact evidence for or against prior reporting of the pulse-timing interaction. If already established, abandon its novelty claim. \textit{ (Not listed by author; not presumed resolved; 3 saved source reference(s))}
\item Validate basal-epidermal exposure and response measurements, noninjurious pulse delivery and stable response-score normalization. If early and late exposures cannot be matched, stop rather than interpret their contrast as a timing effect. \textit{ (Not listed by author; not presumed resolved; 3 saved source reference(s))}
\item Use a separate pilot to justify the minimum-interest margin, estimate between-animal and between-clutch variance, and choose a finite confirmatory sample size for the interaction. Do not select the response module or effect threshold using confirmatory outcomes. \textit{ (Not listed by author; not presumed resolved; 3 saved source reference(s))}
\item Obtain independent randomized molecular and regeneration results with uncertainty. Resolve cell composition, developmental progression, metabolic state and injury-severity alternatives; a pulse rescue alone does not identify the decoder. \textit{ (Not listed by author; not presumed resolved; 3 saved source reference(s))}
\item Animal-level prediction and any claim explaining later regained competence require additional longitudinal measurements and a verified later-competent X. laevis condition, respectively. \textit{ (Not listed by author; not presumed resolved; 3 saved source reference(s))}
\item Audit original GEO metadata, biological units, pairing, pools, batches, phenotypes and file availability. \textit{ (Not listed by author; not presumed resolved; 4 saved source reference(s))}
\item Inspect original and subsequent literature for exact overlap with the pulse-timing interaction. \textit{ (Not listed by author; not presumed resolved; 4 saved source reference(s))}
\item Validate pulse delivery, local exposure, basal-epidermal assays and normalization; separate wound age from post-pulse response latency. \textit{ (Not listed by author; not presumed resolved; 4 saved source reference(s))}
\item Justify the minimum-interest margin and determine confirmatory sample size using separate pilot animals and clutch-level variability. \textit{ (Not listed by author; not presumed resolved; 4 saved source reference(s))}
\item Obtain independent randomized molecular and regeneration results, with uncertainty and composition, metabolic, developmental and injury controls. \textit{ (Not listed by author; not presumed resolved; 4 saved source reference(s))}
\item Require individual measurement\textendash{}outcome links for animal-level prediction and a verified later-competent condition for regained-competence claims. \textit{ (Not listed by author; not presumed resolved; 4 saved source reference(s))}
\item Material and recoverable: the fixed 6-hour endpoint confounds pulse timing with elapsed time since treatment. Stage-dependent response kinetics can produce the proposed interaction without the claimed wound-age-dependent gate. \textit{ (Not listed by author; not presumed resolved; 2 saved source reference(s))}
\item Material and recoverable: neither accession structure nor prior reporting of the timing interaction has been inspected. The mathematical limitation cannot yet be attributed to the actual GEO samples, and novelty cannot be claimed. \textit{ (Not listed by author; not presumed resolved; 2 saved source reference(s))}
\item Material and recoverable: no independent molecular or regeneration measurements support the candidate, and no verified later-competent X. laevis condition tests regained competence. These are missing evidence, not refutations. \textit{ (Not listed by author; not presumed resolved; 2 saved source reference(s))}
\item The proposed timing contrast confounds wound age with response latency. \textit{ (Not listed by author; not presumed resolved; 2 saved source reference(s))}
\end{itemize}
111 further distinct obligations are retained in the archived reconciliation.
\section{Conclusions}
Unlinked stage-specific early and late distributions do not generally identify animal-level temporal coupling. This conditional deduction limits interpretation of a sampling design; it neither establishes the structure of the nominated GEO datasets nor argues against the existence of biological coupling.

Equal-exposure pulses at different wound ages provide a potentially discriminating test of restricted insufficient-signal and delayed-receptivity models. The predicted transition-drive contrast is conditional, and a molecular timing effect must survive matched-latency controls and competing explanations before receiving a mechanistic interpretation.

The investigation remains unresolved as biological discovery. No numerical biological effect, predictive gain, novelty determination, independent confirmation, or explanation of later recovered competence is established by the supplied evidence.

\textit{Claims discussed: }C10, C11, C12
\section*{Evidence and authorship note}
Snapshot SHA-256: 3236a8c66a96d713f28726271169e67c88c8a768f398ab3b9b504fee3cdcc74e
Generated by an AI author from the cited saved research snapshot. This document is not independent verification. No external bibliography was added.
\section*{Validation warnings}
\begin{itemize}
\item Evidence indexing omitted 3390 material saved research field fragment(s). Those fragments were not reviewed for this paper; the interpretation is provisional where the omitted evidence may matter. Consult the frozen source snapshot and omission ledger for full details.
\item The supplied index is a bounded selection, not evidence of exhaustive review of the saved research or literature.
\item Source-linked model arguments and internal agreement are not independent empirical validation or formal verification.
\item The biological premise of unlinked early and late sampling has not been verified for the named accessions; no accession-specific impossibility claim is warranted.
\item The originally proposed fixed-endpoint timing comparison has a consequential latency confound that remains unresolved without an adequate molecular sampling schedule.
\item The proposed pulse model does not establish that ERK is the operative decoder, that its manipulation is selective, or that a change in transition drive produces functional regrowth.
\item Known RIC/ECM involvement and other pathway associations cannot be presented as new findings, and the timing interaction's prior-art status remains unassessed.
\item No empirical effect size, confidence interval, significance result, sample-size adequacy, or predictive performance can be supplied from this evidence.
\item The final structured claim register must contain unique entries for every referenced claim identifier. The present register is incomplete: C01 through C09 are referenced but absent, and C11 and C12 are duplicated.
\item The evidence index explicitly omitted 3390 saved field fragment(s); see the generation coverage ledger and review the complete archived snapshot.
\item Structured paper indexing selected 447 source-linked text fragments from 15 archived sources without AI extraction calls. The complete frozen snapshot remains archived; this bounded index does not constitute review of every saved field.
\item The frozen evidence contains open proof/research obligation(s) requiring reconciliation: 121 item(s) are not listed by the author and are not presumed resolved; 8 author-listed item(s) have no exact frozen-record match. See the obligation reconciliation.
\item The section Biological question and observational scope narrative has no linked claim IDs; check it for unsupported factual statements.
\item The section Why unlinked snapshots cannot identify animal-level coupling narrative has no linked claim IDs; check it for unsupported factual statements.
\item The section A restricted timing-gate hypothesis narrative has no linked claim IDs; check it for unsupported factual statements.
\end{itemize}
\end{document}
